\documentclass[sigconf]{acmart}

%% Packages for LaTeX Tables to render properly
\usepackage{graphicx}
\usepackage{booktabs}
\usepackage{tabularx}
\usepackage{array}
\usepackage{multirow}
\usepackage{longtable}
\usepackage{ltxtable}
\usepackage{geometry}
\usepackage{ragged2e}
\usepackage{wasysym}
\usepackage[table]{xcolor}
\usepackage{makecell}
\usepackage{pifont}
\usepackage{amssymb}
\usepackage{xcolor}

% Custom Column Types -- for Tables
\newcolumntype{C}[1]{>{\centering\arraybackslash}m{#1}}
\newcolumntype{L}[1]{>{\raggedright\arraybackslash}m{#1}}
\newcolumntype{M}[1]{>{\centering\arraybackslash}m{#1}}
\newcolumntype{G}{!{\color{gray!50}\vrule}}

% Multirow alignment -- for Tables
\renewcommand{\multirowsetup}{\centering\arraybackslash}

% Color Setting -- for Tables
\arrayrulecolor[gray]{0.85}


%% ACM Rights management information.
\copyrightyear{2025}
\acmYear{2025}
\setcopyright{acmlicensed}
\acmConference[DocEng '25]{ACM Symposium on Document Engineering 2025}{September 2--5, 2025}{Nottingham, United Kingdom}
\acmBooktitle{ACM Symposium on Document Engineering 2025 (DocEng '25), September 2--5, 2025, Nottingham, United Kingdom}
\acmDOI{10.1145/3704268.3748682}
\acmISBN{979-8-4007-1351-4/2025/09}

%% Start of document
\begin{document}

\title{Document Encryption in Practice: A Comparative Framework and Evaluation}

\author{Isaac Henry Teuscher}
\orcid{0009-0000-2238-1279}
\affiliation{%
  \institution{Brigham Young University}
  \city{Provo}
  \state{Utah}
  \country{USA}}
\email{isaacteuscher@gmail.com}

\author{Benjamin Schooley}
\orcid{0000-0002-6918-9986}
\affiliation{%
  \institution{Brigham Young University}
  \city{Provo}
  \state{Utah}
  \country{USA}}
\email{ben_schooley@byu.edu}

\begin{abstract}
Document files with sensitive information are used across nearly every industry. In recent years, cyberattacks have resulted in millions of sensitive documents being exposed. Although
document encryption methods exist, they are often flawed in terms of usability, security, or deployability. We present a structured framework for evaluating document encryption methods, adapting the usability-deployability-security (“UDS”) model to the document encryption context. We apply this framework to compare current methods, performing a comprehensive evaluation of nine document protection methods, including password-based, passwordless, and cloud-based approaches. Our analysis across 15 design properties highlights the benefits and limitations of current methods. We propose
strategies and design recommendations to address key limitations such as memory-wise effort, granular protection, and shareability.
\end{abstract}

%% ACM CCS Concepts. The code below is generated by the tool at http://dl.acm.org/ccs.cfm.
\begin{CCSXML}
<ccs2012>
   <concept>
       <concept_id>10002978.10003022</concept_id>
       <concept_desc>Security and privacy~Software and application security</concept_desc>
       <concept_significance>500</concept_significance>
       </concept>
 </ccs2012>
\end{CCSXML}

\ccsdesc[500]{Security and privacy~Software and application security}

%% Keywords
\keywords{Document encryption, Secure document engineering, Security and usability}

%% This command processes the author and affiliation and title information and builds the first part of the formatted document.
\maketitle

\section{Introduction}
Electronic documents containing highly sensitive data are nearly ubiquitous, used across a wide range of industries. Examples include health records, legal counsel, financial data, trade secrets, and employment contracts. These sensitive documents are often stored and shared using cloud services like managed file transfer (MFT) applications. Users trust these platforms to keep their data secure.

However, in recent years, several large MFT applications
experienced mass exploit events that compromised millions of confidential documents, including personal health data, government records, and legal files \cite{kondruss_moveit_2023}. The largest incident resulted in the exposure of 95,788,491 individuals’ personal
information \cite{simas_unpacking_2023}. Although MFT applications encrypt files \textit{'at rest'}, all documents are unencrypted and readable by a privileged user session. From 2020 to 2025, five vulnerabilities allowing
remote privileged access to MFTs were discovered and abused, resulting in major data breaches\footnote{Accellion Dec 2020, Serv-U Nov 2021, GoAnywhere Jan 2023, MOVEit May 2023, and Cleo Dec 2024.} \cite{cisa_cl0p_2023}. These incidents underscore ongoing challenges in the state of document encryption. While
several approaches exist—such as password-based, proprietary, and storage provider encryption—each comes with major limitations. Issues of security, usability, and deployability hinder widespread adoption and reduce effectiveness. Improved methods are needed to address these challenges.

In this paper, an analytical framework for measuring
document protection methods is presented consistent with approaches taken to evaluate authentication schemes in related domains \cite{alaca_comparative_2021-1, alaca_comparative_2021, bonneau_quest_2012, borisov_sok_2021}. This framework provides a semi-structured evaluation methodology and benchmark for future document protection methods \cite{bonneau_quest_2012}. We aim to answer the following question: \textit{What usability, deployability, and security limitations exist in current document encryption methods?}

\section{Related Work}
Several research efforts have demonstrated weaknesses with current file security measures, including digital signatures \cite{rohlmann_breaking_2021}, XML
encryption, and PDF encryption \cite{muller_practical_2019}. Work has been done on document integrity and non-repudiation both to improve the integrity of documents \cite{shatnawi_maintaining_2017} and to uncover vulnerabilities in integrity features such as PDF certification signatures \cite{rohlmann_breaking_2021}. This paper focuses on the challenge of maintaining confidentiality instead of integrity or non-repudiation. Little research has been done to directly assess the usability of document encryption methods. However, extensive usability research has been performed on passwords \cite{bonneau_quest_2012, zimmermann_quest_2018}, passwordless technology \cite{alaca_comparative_2021-1}, and encrypted email \cite{borisov_sok_2021} which each present important connections to the usability challenges of document encryption.

\input{benefits-table-acm}

\section{Framework}
The usability-deployability-security (“UDS”) framework was first presented by Bonneau et al. \cite{bonneau_quest_2012} to evaluate password replacement schemes and has subsequently been reused \cite{alaca_comparative_2021-1, alaca_comparative_2021}, expanded \cite{zimmermann_quest_2018}, and adapted \cite{borisov_sok_2021}. The methodology draws on usability evaluation methods such as feature inspection and Nielsen’s heuristic analysis to identify common design properties and develop a taxonomy comparing the benefits offered by each method. The
focus is on empirical and technical metrics, limiting subjectivity. Research studies have adapted the UDS framework to specific contexts such as web authentication \cite{alaca_comparative_2021-1}, single sign-on \cite{alaca_comparative_2021}, and encrypted email \cite{borisov_sok_2021}, tailoring the measured benefits to fit the unique needs of each domain.

The categories and evaluation criteria presented in this work are consistent with how UDS frameworks have been developed and applied in prior literature. The development of this document encryption framework involved a review of related standardized evaluation frameworks, an analysis of technical specifications for each product, and expert consultation—with three information security academic researchers and two cybersecurity professionals—to validate decisions.

Several caveats should be considered when interpreting the evaluation framework. First, the ordering does not imply importance or ranking. Second, the benefits chosen are not all of equal weight. Depending on the use case, different criteria may be more or less important (e.g., what matters for an enterprise may not matter for personal consumer use). The framework does not produce precise quantitative scores but enables structured comparison and benchmarking of new methods against existing solutions. The framework is not set in stone; changes in user familiarity, novel security attacks, and other technological advances may require updates over time. As such, the taxonomy is designed to be adaptable and reusable. Ultimately, the framework illuminates the benefits and limitations of current methods and establishes the design requirements for an ideal method that balances usability, deployability, and security.

\subsection{Benefits}
We define 15 benefits or desirable properties of document
encryption systems, which highlight important evaluation
dimensions in Table 1. We considered including additional benefits used in previous UDS frameworks, such as \textit{Physically-Effortless}, \textit{Efficient-to-Use}, and \textit{Accessible}; however, as these benefits are consistently present across current methods and do not provide comparative value, they were omitted from the framework. When evaluating, we assess the benefits provided in typical \textit{‘in-the-wild’} use cases instead of the idealistic use case, as our focus is on the practical benefits provided. For example, we assume users may recycle passwords or choose simple passwords.

\subsection{Groups}
Three primary authentication types are used today for document encryption: \textbf{password-based} (a known secret is provided for authentication), \textbf{cloud-based} (authentication takes place between a client and server), and \textbf{passwordless} (a cryptographic key is provided for authentication). Each authentication type can be applied at various scopes: individual file level, full disk, and file share. These authentication methods and scopes provide the structure for comparing document encryption methods with the added dimension of license (open-source or proprietary).

\input{comparison-table-acm}

\section{Evaluation}
The evaluation—summarized in Table 2—shows that no
document encryption method today provides all benefits. As such, the choice of which method to use depends on the importance of each benefit and the tradeoffs associated with a particular use case. We summarize key findings by method below.

\subsection{Password-based encryption}
One of the most common forms of document encryption today is password-based file-level encryption, which is built into tools like Microsoft Word® and Adobe Acrobat®. Typically, AES 256-bit encryption is used to protect the entire document with a single password (up to 15 characters) used to derive the encryption key.

A major limitation of this method is that, unlike web-based authentication, a password cannot be changed or recovered if forgotten\footnote{https://support.microsoft.com/en-us/topic/you-cannot-recover-a-document-that-is-protected-with-a-password-if-the-password-is-lost-in-word-65a1f34b-fa3f-665a-4913-7ecbb8e01f5f} violating \textit{U4 Easy-Recovery-from-Loss}. This method also lacks \textit{U1 Memorywise-Effortless} and \textit{U2 Scalable-for-Users}, as users must remember unique passwords for each file. In practice, this leads to weak or reused passwords, undermining the potential \textit{S4 Granular-Protection} benefit of file-level encryption.

Password-based encryption also often fails to provide \textit{S2 Resilient-to-Guessing}. Unlike web authentication, which can throttle login attempts, encrypted files are vulnerable to unthrottled offline brute-force attacks. Researchers have studied how to crack the password of an encrypted PDF using GPUs \cite{danczul_cuteforce_2013} and tuned software \cite{hranicky_experimental_2024}. This research found that nearly any password-based document encryption can be cracked with “little overhead” given enough time and the proper computational resources \cite{danczul_cuteforce_2013}. While some implementations use strong key derivation functions—such as PBKDF2 or Argon2—and high iterations, which make brute-force guessing non-feasible,
dictionary attacks remain effective due to users’ tendency to choose common, simple passwords.

Password-based encryption can also be applied at the full disk level, which does not provide \textit{S4 Granular-Protection} or \textit{U4 Easy-Recovery-from-Loss} and is typically used for personal device protection instead of sharing documents with others.

Given the usability and security limitations of password-based methods, one may wonder why they are still widely used. It is important to note that password-based encryption performs very well in the deployability category. Boneau et al. [3] found that passwords remain a dominant web-based authentication scheme largely because no alternative method provides the same or better levels of deployability benefits. The same holds true in this evaluation of document encryption methods. The familiar password-based encryption of Microsoft Word and PDF documents offers negligible costs (\textit{D1}, \textit{D2}), broad compatibility (\textit{D3}), and wide adoption (\textit{D4}). Perhaps most importantly, it is one of the few methods that provides a level of the \textit{D6 Guest-Friendly} benefit—the ability for a third party to access shared items without requiring provisioning or unique credentials. While the password must be shared, no user account is needed (thus it partially satisfies \textit{D6}). An example illustrating the importance of this benefit is the United States Department of Justice’s instructions for submitting Claims Collection Litigation Reports. These instructions advise users to create a PDF document, encrypt it with a password, and then to “send the Password ... in a separate email for decryption.... Never send the decryption Password in the same email as the encrypted files”\footnote{https://www.justice.gov/jmd/file/789256/dl} . This guidance highlights the lack of usability inherent in password-based document encryption methods but also demonstrates the need for document protection and the reasoning why password-based file encryption remains
widely used in this context: it’s free, it’s interoperable, and it’s simple to use for both the public and government agencies.

\subsection{Cloud-based encryption}
Cloud- or server-based encryption methods (such as MFTs and cloud storage applications) are a second popular document encryption method used today. There are several security and usability benefits that the software-as-a-service (SaaS) model provides. Many of the benefits come from the authentication process functioning at the identity level instead of the individual file level. Often SSO (single sign-on) is used, which can provide the benefit of \textit{S1 Resilient-to-Observation} by enforcing MFA
(multi-factor authentication). Access to documents is not confined to a single device, and account recovery provides effective \textit{U4 Easy-Recovery-from-Loss}. Furthermore, cloud-based models typically provide familiar and effective methods for sharing files
within or outside of an organization. Cloud-based encryption methods are the only schemes to provide both \textit{U5 Seamless-Sharing} and \textit{D6 Guest-Friendly}. Both of which are frequently valued in enterprise and personal use cases. 

A major limitation of cloud-based encryption methods is the lack of \textit{S4 Granular-Protection} and \textit{S3 Resilient-to-Third-Party-Leak}. These security limitations come because the cloud provider handles both encryption and authentication. Once a user’s session
is authenticated, all documents they are authorized to access are effectively decrypted. While a vulnerability in any encryption method can lead to a confidentiality breach, managed cloud-based encryption introduces an additional risk: a breach can occur on the provider’s side even if the customer’s system has no vulnerability.
Additionally, a supply chain incident affecting the provider can impact all customers using the managed service. This scenario occurred five times between 2021 and 2025 as cyber-criminal groups targeting managed file transfer (MFT) applications discovered vulnerabilities and exfiltrated millions of documents from thousands of customers using the same platforms \cite{cisa_cl0p_2023, kondruss_moveit_2023,simas_unpacking_2023}.

Similar to password-based encryption, while there are
limitations of cloud-based encryption, there are other benefits that outweigh the limitations in certain use cases. Again, the importance of deployability and shareability is shown to outweigh security and usability benefits in many real-world scenarios.

\subsection{Passwordless encryption}

Passwordless file-level encryption offers strong security benefits relative to other methods, particularly in terms of \textit{S2 Resilient-to-Guessing}, \textit{S3 Resilient-to-Third-Party-Leak}, and \textit{S4 Granular-Protection}, because it relies on cryptographic private keys rather than identity verification or user-generated passwords.

However, these security benefits come with notable usability and deployability challenges. Sharing passwordless encrypted documents requires secure key exchange, and without integration into widely used platforms, this process often lacks \textit{U5 Seamless-Sharing} and \textit{D6 Guest-Friendly}. Additionally, recovery from lost keys is difficult or impossible, making \textit{U4 Easy-Recovery-from-Loss} a significant limitation.

Today, using passwordless encryption to provide confidentiality of documents remains uncommon in both enterprise and personal domains due to complex deployment and limited interoperability. However, if barriers to adoption, especially around interoperability and sharing, are addressed, passwordless encryption could offer the strongest overall security of any method evaluated in this framework.

\section{Conclusion}
Since password-based encryption was introduced in standard document editing tools, innovation in protecting documents has lagged. While the confidentiality of web traffic has improved over time, the humble document has been left behind in many respects. While digital signatures are commonly used in PDF software to verify the owner or editor of a document (integrity and non-repudiation), protecting the content of a document (confidentiality) often still relies on a simple, fixed password. An open-source, passwordless document encryption method that integrates with widely used identity providers could address many limitations of password-based systems. The design of such a method remains an open problem. This paper proposes a framework to guide the development of secure, interoperable, and user-friendly encrypted document systems and provides a benchmark for comparing current and future methods.

%% Bibliography
  \bibliographystyle{ACM-Reference-Format}
  \bibliography{DocEng25-bibliography}

\end{document}